% Generated by scripts/generate_assets.py; do not edit.
\begin{table}[!htbp]
\centering
\normalsize
\caption{Anualização contábil da compensação de produtividade}\label{tab:tfp_horizons}
\setlength{\tabcolsep}{4pt}
\begin{tabular}{llrrrr}
\toprule
Função & Teto & $A_{\rm req}$ (\%) & 3 anos & 5 anos & 10 anos \\
\midrule
Bilateral & 40h & 1,573 & 0,521 & 0,313 & 0,156 \\
Bilateral & 36h & 8,381 & 2,719 & 1,623 & 0,808 \\
Acima do pico & 40h & 1,534 & 0,509 & 0,305 & 0,152 \\
Acima do pico & 36h & 6,524 & 2,129 & 1,272 & 0,634 \\
\bottomrule
\end{tabular}
\par\smallskip
\begin{minipage}{0.98\linewidth}\normalsize Os ganhos anuais em \% são $g_T=100[(1+A_{\rm req})^{1/T}-1]$, com $A_{\rm req}$ em fração. Base: horas formais limitadas a 44h. A anualização apenas expressa a mesma diferença de nível em outra escala; não é uma trajetória dinâmica estimada. Na PWT 11.0 via FRED, o crescimento anualizado de produtividade do Brasil foi -0,83\% em 1990--2023. Entre janelas de dez anos inteiramente contidas em 1990--2019, a maior taxa foi 0,05\% em 2000--2010. A série serve como referência descritiva de escala e não valida o contrafactual.\end{minipage}
\end{table}
